\documentclass[11pt]{article}

\usepackage[T1]{fontenc}
\usepackage{amsmath,amssymb,amsthm}
\usepackage{hyperref}

\title{Profile Likelihood-Ratio Intervals, the Bartlett Correction,\\
and the Eigenvalue-Spectrum Reference Law}
\author{}
\date{\today}

\newcommand{\E}{\mathbb{E}}
\newcommand{\tr}{\operatorname{tr}}
\newcommand{\Var}{\operatorname{Var}}
\newcommand{\Cov}{\operatorname{Cov}}
\newcommand{\vech}{\operatorname{vech}}
\newcommand{\argmin}{\operatorname*{arg\,min}}
\newcommand{\argmax}{\operatorname*{arg\,max}}
\newcommand{\plim}{\operatorname*{plim}}
\newcommand{\Gnt}{\Gamma_{\mathrm{NT}}}

\theoremstyle{plain}
\newtheorem{theorem}{Theorem}
\newtheorem{lemma}{Lemma}
\newtheorem{corollary}{Corollary}
\theoremstyle{remark}
\newtheorem*{remark}{Remark}

\begin{document}
\maketitle

\section*{Purpose}

Records the theory behind the experiment-43 maximal-reliability interval: the
profile likelihood-ratio (LR) statistic and its inversion into a confidence set,
the Bartlett correction that rescues its small-$N$ calibration under normality,
and what changes under non-normality, where the reference law is a weighted sum
of $\chi^2_1$ variables (the full eigenvalue spectrum, of which Satorra--Bentler
scaling is only the first-moment summary). Three claims, stated here and argued
below.

\begin{enumerate}
\item Test inversion buys transformation-invariance, range-respecting limits,
and asymmetry for free, by anchoring on the likelihood rather than on
$\hat\psi$ (Section~\ref{sec:profile}).
\item The Bartlett correction is a finite-sample ($O(N^{-1})$) fix that is valid
\emph{only when the limit is genuinely $\chi^2$}. It removes the small-$N$
inflation of the LR and, uniquely among the trinity, does so to one extra order
(Section~\ref{sec:bartlett}).
\item Under non-normality the leading-order ($O(1)$) error must be fixed first by
a spectral correction (Section~\ref{sec:spectrum}). A single scalar Bartlett
factor on the mixture is then only a first-moment patch, not a Bartlett
correction in the Barndorff-Nielsen sense; the principled finite-sample-robust
route is to correct (or resample) the eigenvalues, not to rescale a mean
(Section~\ref{sec:bartlett-spectrum}). For a \emph{one}-parameter interval the
spectrum collapses to a single eigenvalue, so there the robust correction is just
a scalar and the distinction is moot; it bites only for $\mathrm{df}>1$ (overall
fit, multi-parameter regions, nested tests).
\end{enumerate}

\section{The profile LR and test-inversion intervals}
\label{sec:profile}

Partition $\theta = (\psi, \lambda)$ with $\psi \in \mathbb{R}$ the interest
parameter (here the maximal reliability $\rho^*$) and $\lambda$ the nuisance
block. In covariance-structure ML the discrepancy $F$ satisfies
$-2\,\ell(\theta)/N = F(\theta) + \text{const}$, and we write
$\hat\theta = \argmin_\theta F$ for the unconstrained fit and
$\tilde\theta(\psi_0) = \argmin_{\theta:\,\psi=\psi_0} F$ for the fit constrained
to $\psi = \psi_0$. The profile discrepancy is
$F_p(\psi) = \min_\lambda F(\psi,\lambda)$, and the profile LR statistic is
%
\begin{equation}
\label{eq:lr}
W(\psi_0) \;=\; N\big[F_p(\psi_0) - F_p(\hat\psi)\big]
          \;=\; N\big[F(\tilde\theta(\psi_0)) - F(\hat\theta)\big],
\end{equation}
%
which is exactly the experiment-43 statistic $\mathrm{LR} = N(F_{\mathrm{con}} -
F_{\mathrm{unc}})$.

\begin{theorem}[Wilks]
\label{thm:wilks}
Under the usual regularity conditions, correct specification, normal-theory ML,
and $\theta$ interior with true interest value $\psi_0$,
$W(\psi_0) \xrightarrow{d} \chi^2_q$ with $q = \dim\psi$ ($=1$ here)
\cite{wilks1938}.
\end{theorem}

The $(1-\alpha)$ confidence set is the acceptance region of the level-$\alpha$ LR
test, swept over the hypothesized value:
%
\begin{equation}
\label{eq:ci}
C_{1-\alpha} \;=\; \big\{\psi_0 : W(\psi_0) \le \chi^2_{q,\,1-\alpha}\big\}.
\end{equation}
%
Coverage of the true value needs only $W$ \emph{at} that value, which is why
experiment-43 measures coverage with one constrained fit per replicate and defers
interval width (the latter needs the root-finding $W(\psi_0) = \chi^2_{q,1-\alpha}$).

Four properties make \eqref{eq:ci} succeed where intervals built around
$\hat\psi$ fail.
\begin{enumerate}
\item[(i)] \emph{Transformation-invariance.} $W$ is invariant under any monotone
reparameterization of $\psi$; a reparameterization maps $C$ to the image set. No
logit-versus-Fisher choice arises.
\item[(ii)] \emph{Range-respecting.} The constrained fit stays in the parameter
space, so $C \subseteq [0,1]$ for a reliability by construction; the interval
cannot leave the natural range.
\item[(iii)] \emph{Asymmetry.} $C$ follows the curvature of $F_p$ and is
asymmetric whenever the profile likelihood is.
\item[(iv)] \emph{No standard error.} Only fits enter, never an estimate of
$\Var(\hat\psi)$, so $C$ is immune to the bias, skew, and boundary pileup of
$\hat\psi$ that defeat Wald, delta, and Cornish--Fisher intervals.
\end{enumerate}

\begin{remark}
$W$, the Wald statistic $W_W = (\hat\psi - \psi_0)^2/\widehat{\Var}(\hat\psi)$,
and the score statistic are asymptotically equivalent to $o_p(1)$ and share the
limit of Theorem~\ref{thm:wilks}. $W_W$ is the quadratic Taylor approximation of
$W$ at $\hat\psi$; it is \emph{not} transformation-invariant and inherits every
finite-sample defect of $\hat\psi$. The signed root
$R(\psi_0) = \operatorname{sign}(\hat\psi - \psi_0)\sqrt{W(\psi_0)} \to N(0,1)$ is
the one-sided handle and the entry point for higher-order ($r^*$) refinements.
In SEM, profile-LR intervals are due to \cite{nealemiller1997} (genetics, the
OpenMx \texttt{mxCI} machinery) and \cite{pekwu2015} (general SEM, with
confidence regions); the R package \texttt{semlbci} implements them.
\end{remark}

\section{The Bartlett correction}
\label{sec:bartlett}

Theorem~\ref{thm:wilks} is leading-order only. The next term is a pure scale
effect:
%
\begin{equation}
\label{eq:bartlett-mean}
\E[W(\psi_0)] \;=\; q\Big(1 + \tfrac{b}{N} + O(N^{-2})\Big),
\end{equation}
%
with $b$ a model-dependent constant. At small $N$ the mean exceeds $q$, $W$ is
stochastically too large, the test over-rejects, and $C$ under-covers. In
experiment-43 this is stark: $\E[W] \approx 3.8$ against $q = 1$ at $N = 50$ for
the general-factor coefficient, decaying to $\approx 1$ by $N = 200$.

\begin{theorem}[Bartlett-correctability]
\label{thm:bartlett}
Let $c = 1 + b/N = \E[W]/q$ and $W' = W/c$. Then $W' \xrightarrow{d} \chi^2_q$
with relative error $O(N^{-2})$, one order better than for $W$
\cite{bartlett1937,lawley1956}. The single mean-matching rescaling that sets
$\E[W'] = q$ simultaneously corrects the higher cumulants to the order needed; the
level error after adjustment is $O(N^{-2})$ \cite{barndorffnielsencox1984,
barndorffnielsenhall1988}. This is special to the LR: the Wald and score
statistics are not Bartlett-correctable by a scalar.
\end{theorem}

The correction is the rescaling $W \mapsto W/c$ in \eqref{eq:ci}, equivalently a
rescaled threshold $c\,\chi^2_{q,1-\alpha}$. The factor $c = \E[W]/q$ can be
obtained three ways:
\begin{itemize}
\item analytically, from Lawley's general expansion of \eqref{eq:bartlett-mean}
(model-specific and algebraically heavy);
\item by parametric bootstrap, simulating $W$ under the fitted null and taking the
sample mean;
\item empirically/oracle, the Monte-Carlo mean of $W$ at the true value, as in
experiment-43 (an oracle, not an operational estimator).
\end{itemize}

\begin{remark}[SEM lineage]
The covariance-structure analogue is old. \cite{bartlett1950} gave the
factor-analysis $\chi^2$ multiplier; \cite{swain1975} the Swain correction (a
Bartlett-type scalar for the ML $\chi^2$, evaluated by \cite{nevitthancock2004,
herzogboomsmareinecke2007}); \cite{yuantianyanagihara2015} an empirical
correction tuned for many variables. All target the same $O(N^{-1})$
over-rejection of the normal-theory $\chi^2$. The standing assumption everywhere
is that the limit \emph{is} $\chi^2$: a Bartlett correction fixes the
finite-sample approach to that limit and presupposes it exists.
\end{remark}

\section{The reference law under non-normality: the full spectrum}
\label{sec:spectrum}

Under non-normality (or misspecification) the normal-theory statistic no longer
converges to $\chi^2$. Let $\Gamma = \Cov\!\big(\sqrt{N}\,\vech(S - \Sigma)\big)$
be the asymptotic covariance of the sample moments, $\Delta = \partial\,
\vech\Sigma(\theta)/\partial\theta'$ the model Jacobian, $V$ the fit weight
(normal-theory $V = \Gnt^{-1}$ for ML), and
%
\begin{equation}
\label{eq:U}
U \;=\; V - V\Delta(\Delta'V\Delta)^{-1}\Delta'V
\end{equation}
%
the residual-weight (projection) matrix.

\begin{theorem}[Spectral limit]
\label{thm:spectrum}
Under correct structural specification and finite fourth moments,
%
\begin{equation}
\label{eq:mixture}
T \;=\; N\hat F \;\xrightarrow{d}\; \sum_{j=1}^{r} \lambda_j\, z_j^2,
\qquad z_j \stackrel{\text{iid}}{\sim} N(0,1),
\end{equation}
%
where $\lambda_1,\dots,\lambda_r$ are the $r = \mathrm{df}$ nonzero eigenvalues of
$U\Gamma$ \cite{shapiro1983,satorra1989,satorrabentler1994}. Under multivariate
normality $\Gamma = \Gnt$ and $U\Gnt$ is idempotent of rank $r$, so every
$\lambda_j = 1$ and $T \to \chi^2_r$. Non-normality moves the eigenvalues off $1$.
\end{theorem}

So the statistic is a weighted sum of independent $\chi^2_1$, not a $\chi^2$. The
established corrections form a hierarchy by how much of the spectrum
$\{\lambda_j\}$ they use:
\begin{itemize}
\item \emph{Mean} (Satorra--Bentler $T_M$): scale by
$c_1 = \tr(U\Gamma)/r = \bar\lambda$ and refer to $\chi^2_r$. First moment only;
the variance is wrong unless the $\lambda_j$ are equal.
\item \emph{Mean-and-variance} (Satorra--Bentler adjusted / Satterthwaite
$T_{MV}$): match two moments by scaling and setting fractional
$d = (\sum_j\lambda_j)^2/\sum_j\lambda_j^2$; refer to $\chi^2_d$.
\item \emph{Scaled-shifted} ($T_{MV2}$): match two moments by an affine $aT+b$ at
fixed df \cite{asparouhovmuthen2010}.
\item \emph{Higher moments}: third-moment-adjusted variants extend the same idea.
\item \emph{Full spectrum}: use all $\lambda_j$ and the exact distribution of
$\sum_j\lambda_j z_j^2$ by numerical inversion of its characteristic function
\cite{imhof1961,davies1980}, with no moment truncation.
\end{itemize}

\begin{remark}[Why the full spectrum is hard: the FMG program]
The exact tail of \eqref{eq:mixture} needs the population $\lambda_j$, but the
sample eigenvalues of $\hat U\hat\Gamma$ are badly biased and unstable, worst for
the small eigenvalues and exactly where $N/\mathrm{df}$ is small; plugging them
into Imhof is then worse than moment-matching. Eigenvalue block averaging (EBA)
bins the ordered sample eigenvalues and replaces each by its block mean, trading
variance for bias \cite{foldnesgronneberg2018}; penalized EBA (pEBA / pOLS)
shrinks the spectrum toward a model-based trend and holds the nominal rate where
$T_M$ and $T_{MV}$ drift under non-normality \cite{foldnesmossgronneberg2025},
with the nested-model extension in \cite{foldnesmossgronneberg2026}. magmaan
implements this family (experiments 17--18) to $\sim$$3\times10^{-6}$ parity with
\texttt{semTests}. Satorra--Bentler is, in this language, the rank-collapsed
first-moment summary of the same spectrum the full method resolves.
\end{remark}

\section{A Bartlett correction for the spectral statistic?}
\label{sec:bartlett-spectrum}

The two corrections attack different orders and are orthogonal. The spectral
correction fixes the leading-order ($O(1)$) reference law: it makes the
\emph{asymptotic} distribution right under non-normality. The Bartlett correction
fixes the next-order ($O(N^{-1})$) approach to a $\chi^2$ limit. Under non-normal,
small-$N$ data \emph{both} are wrong, and one needs both. Three things
``Bartlett the spectrum'' can mean.

\paragraph{(a) A first-moment finite-sample patch (exists; recommended).}
Apply a Bartlett-type scalar to an \emph{already-robust} base statistic:
$T_{MB} = $ Bartlett-corrected Satorra--Bentler. \cite{jobstauerswaldmoshagen2022}
compare the family and recommend exactly this ordering, mean-correct (SB) first,
then a Bartlett scalar on the scaled statistic, reporting that $T_{MB}$ is
satisfactory across sources of non-normality where bare SB and the moment-matched
variants drift at small $N$. This is the operational answer and the direct sibling
of experiment-43, with the SB mean replacing the normal-theory mean in the factor.

\paragraph{(b) A true Bartlett correction of the mixture (does not exist as a
scalar).} The Barndorff-Nielsen property of Theorem~\ref{thm:bartlett}, that one
scalar fixes the whole distribution to $O(N^{-2})$, is a property of the LR with a
$\chi^2$ limit. For a general weighted-$\chi^2$ limit there is no single scalar
whose mean-match also corrects the higher cumulants: the finite-sample error lives
in the whole eigenvalue vector (each $\lambda_j$ carries its own $O(N^{-1})$ bias),
not in one mean. Properly ``Bartlett-correcting the spectrum'' therefore means
finite-sample-correcting the eigenvalues (or $\Gamma$ and $U$) before Imhof, which
is precisely the bias problem EBA/pEBA already attack by shrinkage; pEBA is, in
effect, the spectral analogue of a finite-sample correction, done by penalization
rather than by an $O(N^{-1})$ expansion.

\paragraph{(c) Resample the reference law (the clean robust-plus-finite-sample
route).} A model-based or nonparametric bootstrap \cite{bollenstine1992} estimates
the finite-sample null distribution of $T$ directly, absorbing non-normality
(leading order) and the $O(N^{-1})$ term at once; the bootstrap mean of $T$ is a
Bartlett factor that is robust by construction. This is why the empirical $c$ of
experiment-43 is more than a normal-theory Bartlett factor: re-estimated on
non-normal data it would absorb the leading-order mis-scaling too. A double
bootstrap refines the calibration further.

\begin{remark}[The df${=}1$ collapse]
For a one-parameter interval ($q = 1$), the constrained-LR limit under
non-normality is a \emph{single} scaled variable $\lambda\, z^2$: the spectrum of
the relevant $q\times q$ matrix has one element. There the robust correction is
just the scalar scaling $\lambda$ (a Satorra-2000 scaled difference test), the
``full spectrum'' has nothing extra to offer, and $T_{MB}$ of path~(a) is the
complete recipe. The spectral richness of Section~\ref{sec:spectrum} matters only
for $\mathrm{df}>1$: overall goodness-of-fit, multi-parameter confidence regions,
and nested tests with more than one restriction.
\end{remark}

\paragraph{Recommendation for a built surface (experiment 43 /
\texttt{measures::frontier}).}
\begin{itemize}
\item Normal data: profile-LR interval \eqref{eq:ci} with a Bartlett factor
(analytic or parametric-bootstrap $c$). A genuine Bartlett correction, $O(N^{-2})$.
\item Non-normal data: do \emph{not} Bartlett the normal-theory LR. Either
(i)~scale by the robust $\lambda$ (Satorra-2000) then Bartlett ($T_{MB}$,
path~(a)), or (ii)~bootstrap the profile-LR null (path~(c)). Both fix leading
order first.
\item The eigenvalue machinery already in magmaan (Imhof tail, EBA/pEBA) is the
spectral reference; the open piece is wiring a spectral (not $\chi^2_1$) threshold
into the profile-LR inversion, the non-normal sibling that the experiment-43
caveat defers.
\end{itemize}

\begin{thebibliography}{99}

\bibitem{wilks1938} Wilks, S.~S. (1938). The large-sample distribution of the
likelihood ratio for testing composite hypotheses. \emph{Annals of Mathematical
Statistics}, 9(1), 60--62.
\href{https://doi.org/10.1214/aoms/1177732360}{doi:10.1214/aoms/1177732360}.

\bibitem{nealemiller1997} Neale, M.~C., \& Miller, M.~B. (1997). The use of
likelihood-based confidence intervals in genetic models. \emph{Behavior
Genetics}, 27(2), 113--120.
\href{https://doi.org/10.1023/A:1025681223921}{doi:10.1023/A:1025681223921}.

\bibitem{pekwu2015} Pek, J., \& Wu, H. (2015). Profile likelihood-based
confidence intervals and regions for structural equation models.
\emph{Psychometrika}, 80(4), 1123--1145.
\href{https://doi.org/10.1007/s11336-015-9461-1}{doi:10.1007/s11336-015-9461-1}.

\bibitem{bartlett1937} Bartlett, M.~S. (1937). Properties of sufficiency and
statistical tests. \emph{Proceedings of the Royal Society of London A}, 160(901),
268--282.
\href{https://doi.org/10.1098/rspa.1937.0109}{doi:10.1098/rspa.1937.0109}.

\bibitem{lawley1956} Lawley, D.~N. (1956). A general method for approximating to
the distribution of likelihood ratio criteria. \emph{Biometrika}, 43(3/4),
295--303.
\href{https://doi.org/10.1093/biomet/43.3-4.295}{doi:10.1093/biomet/43.3-4.295}.

\bibitem{barndorffnielsencox1984} Barndorff-Nielsen, O.~E., \& Cox, D.~R. (1984).
Bartlett adjustments to the likelihood ratio statistic and the distribution of the
maximum likelihood estimator. \emph{Journal of the Royal Statistical Society B},
46(3), 483--495.

\bibitem{barndorffnielsenhall1988} Barndorff-Nielsen, O.~E., \& Hall, P. (1988).
On the level-error after Bartlett adjustment of the likelihood ratio statistic.
\emph{Biometrika}, 75(2), 374--378.
\href{https://doi.org/10.1093/biomet/75.2.374}{doi:10.1093/biomet/75.2.374}.

\bibitem{bartlett1950} Bartlett, M.~S. (1950). Tests of significance in factor
analysis. \emph{British Journal of Psychology, Statistical Section}, 3(2),
77--85.
\href{https://doi.org/10.1111/j.2044-8317.1950.tb00285.x}{doi:10.1111/j.2044-8317.1950.tb00285.x}.

\bibitem{swain1975} Swain, A.~J. (1975). \emph{Analysis of parametric structures
for variance matrices}. PhD thesis, University of Adelaide.

\bibitem{nevitthancock2004} Nevitt, J., \& Hancock, G.~R. (2004). Evaluating
small sample approaches for model test statistics in structural equation modeling.
\emph{Multivariate Behavioral Research}, 39(3), 439--478.
\href{https://doi.org/10.1207/S15327906MBR3903_3}{doi:10.1207/S15327906MBR3903\_3}.

\bibitem{herzogboomsmareinecke2007} Herzog, W., Boomsma, A., \& Reinecke, S.
(2007). The model-size effect on traditional and modified tests of covariance
structures. \emph{Structural Equation Modeling}, 14(3), 361--390.
\href{https://doi.org/10.1080/10705510701301602}{doi:10.1080/10705510701301602}.

\bibitem{yuantianyanagihara2015} Yuan, K.-H., Tian, Y., \& Yanagihara, H. (2015).
Empirical correction to the likelihood ratio statistic for structural equation
modeling with many variables. \emph{Psychometrika}, 80(2), 379--405.
\href{https://doi.org/10.1007/s11336-013-9386-5}{doi:10.1007/s11336-013-9386-5}.

\bibitem{shapiro1983} Shapiro, A. (1983). Asymptotic distribution theory in the
analysis of covariance structures (a unified approach). \emph{South African
Statistical Journal}, 17(1), 33--81.

\bibitem{satorra1989} Satorra, A. (1989). Alternative test criteria in covariance
structure analysis: a unified approach. \emph{Psychometrika}, 54(1), 131--151.
\href{https://doi.org/10.1007/BF02294453}{doi:10.1007/BF02294453}.

\bibitem{satorrabentler1994} Satorra, A., \& Bentler, P.~M. (1994). Corrections
to test statistics and standard errors in covariance structure analysis. In A.
von Eye \& C.~C. Clogg (Eds.), \emph{Latent Variables Analysis} (pp.\ 399--419).
Sage.

\bibitem{asparouhovmuthen2010} Asparouhov, T., \& Muth\'en, B. (2010).
\emph{Simple second order chi-square correction}. Mplus Technical Appendix.
\href{https://www.statmodel.com/download/WLSMV_new_chi21.pdf}{statmodel.com}.

\bibitem{imhof1961} Imhof, J.~P. (1961). Computing the distribution of quadratic
forms in normal variables. \emph{Biometrika}, 48(3/4), 419--426.
\href{https://doi.org/10.2307/2332763}{doi:10.2307/2332763}.

\bibitem{davies1980} Davies, R.~B. (1980). Algorithm AS 155: the distribution of a
linear combination of $\chi^2$ random variables. \emph{Applied Statistics},
29(3), 323--333.
\href{https://doi.org/10.2307/2346911}{doi:10.2307/2346911}.

\bibitem{foldnesgronneberg2018} Foldnes, N., \& Gr\o nneberg, S. (2018).
Approximating test statistics using eigenvalue block averaging. \emph{Structural
Equation Modeling}, 25(1), 101--114.
\href{https://doi.org/10.1080/10705511.2017.1373021}{doi:10.1080/10705511.2017.1373021}.

\bibitem{foldnesmossgronneberg2025} Foldnes, N., Moss, J., \& Gr\o nneberg, S.
(2025). Improved goodness of fit procedures for structural equation models.
\emph{Structural Equation Modeling}, 32(1).
\href{https://doi.org/10.1080/10705511.2024.2372028}{doi:10.1080/10705511.2024.2372028}.

\bibitem{foldnesmossgronneberg2026} Foldnes, N., Moss, J., \& Gr\o nneberg, S.
(2026). Penalized eigenvalue block averaging: extension to nested model comparison
and Monte Carlo evaluations. \emph{Behavior Research Methods}.
\href{https://doi.org/10.3758/s13428-026-02968-4}{doi:10.3758/s13428-026-02968-4}.

\bibitem{jobstauerswaldmoshagen2022} Jobst, L.~J., Auerswald, M., \& Moshagen, M.
(2022). The effect of latent and error non-normality on corrections to the test
statistic in structural equation modeling. \emph{Behavior Research Methods},
54(5), 2351--2363.
\href{https://doi.org/10.3758/s13428-021-01729-9}{doi:10.3758/s13428-021-01729-9}.

\bibitem{bollenstine1992} Bollen, K.~A., \& Stine, R.~A. (1992). Bootstrapping
goodness-of-fit measures in structural equation models. \emph{Sociological Methods
\& Research}, 21(2), 205--229.
\href{https://doi.org/10.1177/0049124192021002004}{doi:10.1177/0049124192021002004}.

\end{thebibliography}

\end{document}
